\chapter{Guía de instalación y entorno reproducible}
\label{ap:instalacion}

Este apéndice concentra los requisitos de los laboratorios. Las versiones cambian; por eso se registran la versión instalada y la fuente oficial, en lugar de copiar comandos de terceros sin revisión. La ruta recomendada es GNU/Linux. macOS funciona con una máquina de Podman y Windows con WSL2 más una distribución Linux.

\section{Capas del entorno}

\begin{table}[htbp]
\centering
\caption{Herramientas por capacidad.}
\begin{tabularx}{\textwidth}{p{3cm}p{3.7cm}X}
\toprule
\textbf{Capacidad} & \textbf{Herramienta libre} & \textbf{Uso en el curso}\\
\midrule
Código y automatización & Git, Make, Python & Versionado, tareas repetibles y aplicación.\\
Editor y diagramas & VSCodium, PlantUML, Graphviz & Código, C4/UML como texto y exportación.\\
Aplicación y datos & FastAPI, PostgreSQL, RabbitMQ & API, persistencia y eventos.\\
Contenedores & Podman, Buildah & Construcción y ejecución OCI sin demonio central.\\
Forja y CI & Forgejo, Forgejo Runner & Revisión, incidencias y canalizaciones.\\
Orquestación & kubectl, kind, Kubernetes & Despliegue local y reconciliación.\\
Observabilidad & OpenTelemetry, Prometheus, Grafana, Jaeger & Métricas, paneles y trazas.\\
Seguridad & Gitleaks, Syft, Trivy, OWASP ZAP & Secretos, SBOM, vulnerabilidades y DAST.\\
IA local & Ollama, Continue & Generación asistida bajo control del entorno.\\
Documento & TeX Live & Compilación de este libro y entregas técnicas.\\
\bottomrule
\end{tabularx}
\end{table}

\section{Perfil mínimo y perfil ampliado}

El perfil mínimo requiere procesador de 64 bits, 8 GiB de RAM y unos 20 GiB libres para repositorios, imágenes y datos. Permite ejecutar la API, base de datos y pruebas por separado. El perfil ampliado, con 16 GiB o más, permite kind, la pila de observabilidad y un modelo local pequeño. Estas cifras son orientativas: la arquitectura, la cuantización del modelo y los servicios simultáneos cambian el consumo.

Los laboratorios se pueden distribuir entre integrantes o ejecutar por etapas. La carencia de GPU no impide el curso; la IA local funciona más lentamente en CPU y puede reemplazarse por una instancia institucional autorizada. Ninguna evaluación exige pagar una suscripción.

\section{GNU/Linux Debian o Ubuntu}

Revise cada paquete antes de aceptar la instalación. Los nombres varían entre versiones de distribución.

\begin{lstlisting}[language=bash]
sudo apt update
sudo apt install git make python3 python3-venv python3-pip \
  podman podman-compose plantuml graphviz curl jq
git --version
python3 --version
podman --version
plantuml -version
\end{lstlisting}

Las herramientas que no estén en el repositorio se instalan desde su publicación oficial, comprobando firma o suma cuando se proporcione. No canalice directamente un script remoto al intérprete en un equipo institucional sin descargarlo, inspeccionarlo y verificar su procedencia.

\section{macOS}

Homebrew es software libre, pero su instalación es una decisión previa del usuario. Si ya está autorizado:

\begin{lstlisting}[language=bash]
brew install git make python podman podman-compose plantuml graphviz jq
podman machine init
podman machine start
podman info
\end{lstlisting}

La VM de Podman necesita CPU, memoria y disco. Ajuste esos recursos al crearla si va a ejecutar kind u observabilidad. Las rutas montadas deben estar compartidas con la máquina.

\section{Windows con WSL2}

Active WSL2 mediante la documentación oficial de Microsoft e instale Ubuntu o Debian. Después siga la sección GNU/Linux dentro de la terminal WSL. Mantenga el repositorio en el sistema de archivos Linux, por ejemplo \texttt{/home/usuario/proyecto}, para evitar problemas de permisos y rendimiento. VSCodium puede conectarse a WSL mediante extensiones compatibles; también puede editar desde herramientas Linux.

Podman puede instalarse dentro de WSL o como Podman Desktop. Elija una ruta y documente dónde se ejecuta el motor para no confundir \texttt{localhost}, volúmenes y redes.

\section{Entorno Python reproducible}

\begin{lstlisting}[language=bash]
python3 -m venv .venv
source .venv/bin/activate
python -m pip install --upgrade pip
python -m pip install -r requirements-dev.txt
python -m pip freeze
\end{lstlisting}

El archivo bloqueado debe registrar versiones y, en una entrega madura, hashes. Actualizar dependencias es una tarea explícita: se revisan cambios, licencia, vulnerabilidades y pruebas antes de confirmar el nuevo bloqueo.

\section{Comprobación integral}

\begin{lstlisting}[language=bash]
git --version
python --version
podman info
podman run --rm docker.io/library/alpine:3.22 echo contenedor-ok
kubectl version --client
kind version
trivy --version
syft version
gitleaks version
ollama --version
\end{lstlisting}

No todos los comandos son necesarios desde el primer capítulo. Un resultado ``comando no encontrado'' indica una instalación pendiente, no que deba instalarse sin revisar. Registre en \texttt{docs/entorno.md}: sistema, arquitectura, RAM, versiones, fecha y desviaciones.

\section{Diagnóstico seguro}

\begin{enumerate}
  \item Lea el mensaje completo y determine si el fallo es versión, permiso, red, puerto o recurso.
  \item Consulte la documentación oficial de la versión instalada.
  \item Reduzca el caso: compruebe primero una herramienta, luego su integración.
  \item No resuelva permisos usando \texttt{sudo} de forma indiscriminada ni cambie a usuario \texttt{root} dentro de la imagen.
  \item No desactive TLS, firmas o escaneos para ``hacer funcionar'' una descarga.
  \item Registre causa, corrección y prueba de que el arreglo funciona.
\end{enumerate}

\begin{validacion}
La instalación está terminada cuando un segundo integrante, usando \texttt{docs/entorno.md}, puede crear el entorno, ejecutar pruebas y levantar el servicio sin recibir credenciales personales. Una captura no reemplaza el comando reproducible ni su salida esperada.
\end{validacion}
